\section{Conclusion}
In this work, we introduced ResearchAgent, a system designed to assist researchers by generating research ideas, which encompass problem identification, method development, and experiment design. Inspired by the human process of ideation, our approach conducts broad and deep literature reviews, integrates knowledge across domains to foster idea cross-pollination, and employs a community of reviewing agents to iteratively refine the generated ideas. Our evaluations, both human and model-based, demonstrated that ResearchAgent produces ideas that are more creative, valid, and clear compared to baselines. While this initial foray shows promising results, multiple challenges remain to operationalize ResearchAgent in real-world research settings. Practical considerations include scaling the knowledge store to encompass diverse research domains, and keeping it current with the latest publications, through which the system can become adaptable even to emerging fields.


\section*{Limitations}
ResearchAgent has some limitations that we hope to address in future work.

First, recall that we built the entity-centric knowledge store to propose beneficial entities during idea generation; however this store is constructed by extracting entities from the titles and abstracts of a limited number of publications (due to the costs of processing them) thereby precluding a large number of other entities and their interconnectedness. 

In addition, the number of entities that we obtain from the BLINK entity linker~\cite{blink} amounts to 3 per paper on average, indicating limited coverage (it is an open-domain linker after all), although it does exhibit generally strong understanding of scientific contexts, as demonstrated by the improvement achieved by the inclusion of its predictions (See Figures~\ref{fig:main} and~\ref{fig:pairwise}, and Table~\ref{tab:examples}). 

Furthermore, since our ResearchAgent is powered by LLMs, similar to any other approaches based on LLMs, it may hallucinate the generated research ideas. While our proposed ResearchAgent can partially mitigate this problem by augmenting LLMs with additional elements, such as references to the target paper and greater entity-centric knowledge, which help ground the generation process in more accurate and relevant information, validating these generated research ideas with experiments is essential to truly accelerate scientific research.

Moreover, while our iterative refinement process with ReviewingAgents demonstrates promising results, it has inherent limitations in scope. Although we employed diverse perspectives by utilizing 15 ReviewingAgents to evaluate three different ideas (problem, method, and experimental design) with five specific criteria for each, this approach may not fully capture the broad range of potential perspectives and criteria necessary for comprehensive evaluation across all different research domains. As acknowledged in the paper, our selection of criteria was informed by their presumed importance, but conducting an exhaustive exploration of all possible criteria over diverse domains is beyond the scope of this work (given the complexity of categorizing and balancing all relevant factors and perspectives). However, we believe the potential of our modular approach allows for customizing and aligning updated or even new criteria to any specific target domain with novel applications, and we leave further expanding them as future work.

Lastly, our ResearchAgent may be less suited for generating ideas in certain domains, such as theoretical sciences, where mathematical reasoning and proof generation play a central role. However, its flexibility allows for customization through specific instructions, enabling the integration of reasoning-based models and techniques to cater to the needs of theoretical research. For instance, in theoretical mathematics, we can instruct (reasoning-based) LLMs to focus on generating proofs or methods and omit experimental design steps that are less relevant. This as an exciting area for future work, where specialized techniques tailored to each domain could be included to broaden its applicability.

\section*{Ethics Statement}
We are aware that the ResearchAgent may have the potential to be misused for nefarious purposes, such as generating research ideas about new explosives, malicious software, and invasive surveillance tools. Notably, this vulnerability is not unique to our approach but a common challenge faced by existing LLMs that possess significant creative and reasoning capabilities, occasionally generating content that may be deemed undesirable. Consequently, it underscores the necessity to enhance the robustness and safety of LLMs more broadly. 

Also, we recognize the risks of unintentional plagiarism associated with using ResearchAgent, where the system might generate ideas that closely mirror existing research due to the regurgitation of training data. While mitigation strategies, such as integrating access to a comprehensive knowledge base to inform users of prior work, can be employed, we understand that building and maintaining such a resource is inherently complex and may not fully eliminate the risk. To further reduce the possibility of plagiarism, recording and tracking all generated ideas could help identify similarities and guide the model to avoid repetition, though this approach would necessitate explicit user consent.

\section*{Acknowledgements}
This work was supported by the National Research Foundation of Korea (NRF) grant funded by the Korea government (MSIT) (No. RS-2023-00256259), the grant of the Korea Machine Learning Ledger Orchestration for Drug Discovery Project (K-MELLODDY) funded by the Ministry of Health \& Welfare and Ministry of Science and ICT, Republic of Korea (grant number: RS-2024-12345678), the Artificial intelligence industrial convergence cluster development project funded by the Ministry of Science and ICT (MSIT, Korea) \& Gwangju Metropolitan City, and the Institute for Information \& communications Technology Planning \& Evaluation (IITP) grant funded by the Korea government (MSIT) (RS-2019-II190075, Artificial Intelligence Graduate School Program (KAIST)).

